\documentclass[11pt]{article}
\usepackage[margin=1in]{geometry}
\usepackage{amsmath,amssymb,amsthm,mathtools}
\usepackage[T1]{fontenc}
\usepackage[utf8]{inputenc}
\usepackage{enumitem}
\usepackage{booktabs,longtable,array}
\usepackage{hyperref}
\usepackage{xcolor}
\usepackage{microtype}
\hypersetup{colorlinks=true, linkcolor=blue!50!black, citecolor=blue!50!black, urlcolor=blue!50!black}

\newtheorem{theorem}{Theorem}[section]
\newtheorem{proposition}[theorem]{Proposition}
\newtheorem{lemma}[theorem]{Lemma}
\newtheorem{corollary}[theorem]{Corollary}
\newtheorem{claim}[theorem]{Claim}
\theoremstyle{definition}
\newtheorem{definition}[theorem]{Definition}
\newtheorem{assumption}[theorem]{Assumption}
\theoremstyle{remark}
\newtheorem{remark}[theorem]{Remark}

\newcommand{\KL}{\mathrm{KL}}
\newcommand{\Ent}{\mathrm{Ent}}
\newcommand{\Prob}{\mathsf{P}}
\newcommand{\E}{\mathbb{E}}
\newcommand{\R}{\mathbb{R}}
\newcommand{\N}{\mathbb{N}}
\newcommand{\Pcal}{\mathcal{P}}
\newcommand{\Mcal}{\mathcal{M}}
\newcommand{\Qcal}{\mathcal{Q}}
\newcommand{\Zcal}{\mathcal{Z}}
\newcommand{\Fcal}{\mathcal{F}}
\newcommand{\Gcal}{\mathcal{G}}
\newcommand{\W}{\mathcal{W}}
\newcommand{\one}{\mathbf{1}}
\newcommand{\dd}{\,d}
\newcommand{\supp}{\operatorname{supp}}
\newcommand{\argmin}{\operatorname*{arg\,min}}
\newcommand{\argmax}{\operatorname*{arg\,max}}
\newcommand{\esssup}{\operatorname*{ess\,sup}}
\newcommand{\ri}{\operatorname{ri}}
\newcommand{\dom}{\operatorname{dom}}
\newcommand{\Law}{\operatorname{Law}}
\newcommand{\Var}{\operatorname{Var}}
\newcommand{\Cov}{\operatorname{Cov}}
\newcommand{\diag}{\operatorname{diag}}
\newcommand{\Id}{\mathrm{Id}}
\newcommand{\Lip}{\operatorname{Lip}}
\newcommand{\osc}{\operatorname{osc}}
\newcommand{\diam}{\operatorname{diam}}
\newcommand{\proj}{\operatorname{proj}}
\newcommand{\e}{\mathrm{e}}

\newenvironment{argument}{\paragraph{Argument.}\begin{enumerate}[leftmargin=2em]}{\end{enumerate}}
\newenvironment{proofoutline}{\paragraph{Proof architecture.}\begin{enumerate}[leftmargin=2em]}{\end{enumerate}}


% Source-faithfulness apparatus used by the paper reconstructions.
\newcommand{\sourceanchor}[1]{\par\smallskip\noindent\textbf{Source anchor.} #1\par\smallskip}
\newcommand{\sourcecomment}[1]{\begin{remark}#1\end{remark}}
\newenvironment{sourceguide}
  {\begin{description}[style=nextline,leftmargin=3.2cm,labelwidth=3cm]}
  {\end{description}}

% Backward-compatible environments used by some of the older reconstructions.
\newenvironment{distilled}[1][]{\begin{theorem}[#1]}{\end{theorem}}
\newenvironment{distillation}{\begin{enumerate}[leftmargin=2em]}{\end{enumerate}}

\newenvironment{commentarynotes}{\begin{remark}\begin{enumerate}[leftmargin=2em]}{\end{enumerate}\end{remark}}

\title{Blanchet--Murthy, \emph{Quantifying Distributional Model Risk via Optimal Transport}:\\
Theorem 1 (Strong Duality), in full}
\author{Distilled literature note --- the baseline our Lean proof is compared against}
\date{Source: \texttt{papers/BlanchetMurthy2019-1604.01446.pdf}, \S2.2--2.3 (pp.~4--7), \S4 (pp.~17--26)}
\begin{document}
\maketitle

\section*{Purpose of this note}
The mathematics \textbf{as it appears in \cite{BlanchetMurthy2019}}.  It is the independent baseline
for \texttt{ForMathlib.OT.dualValue\_le\_droValue} (which discharges the Wasserstein \texttt{hge}),
\texttt{BlanchetMurthy2019.wdro\_strong\_duality} and \texttt{GaoKleywegt2023.strong\_duality\_thm1}.
Our Lean proof of the same theorem takes a \emph{completely different route} (\S\ref{sec:diff}); this
note records theirs so the two can be diffed.

\section{Setting (paper's own)}

$S$ is a Polish space, $\mu\in\Pcal(S)$ the baseline measure, $\delta>0$ the radius.

\paragraph{Assumption (A1).}  $c:S\times S\to\R_+$ is nonnegative, lower semicontinuous, and
\[
  \boxed{\;c(x,y)=0 \iff x=y.\;}
\]

\paragraph{Assumption (A2).}  $f\in L^1(d\mu)$ is upper semicontinuous.

\paragraph{Optimal transport cost.}  $d_c(\mu_1,\mu_2):=\inf\{\int c\dd\pi : \pi\in\Pi(\mu_1,\mu_2)\}$.
For nonnegative lower semicontinuous $c$ the infimum \textbf{is attained} (Villani (2008), Thm 4.1).

\paragraph{Primal and dual.}  With
$\Phi_{\mu,\delta}:=\{\pi\in\Pcal(S\times S): \text{first marginal }\mu,\ \int c\dd\pi\le\delta\}$ and
$\Lambda_{c,f}:=\{(\lambda,\varphi):\lambda\ge0,\ \varphi\in mU(S;\R),\ f(y)-\lambda c(x,y)\le\varphi(x)\ \forall x,y\}$,
\[
  I:=\sup\Big\{\textstyle\int f(y)\dd\pi(x,y) : \pi\in\Phi_{\mu,\delta}\Big\},
  \qquad
  J:=\inf\Big\{\lambda\delta+\textstyle\int\varphi\dd\mu : (\lambda,\varphi)\in\Lambda_{c,f}\Big\}.
\]

\section{Theorem 1, verbatim}
\begin{quote}
Under (A1) and (A2):
\begin{enumerate}[label=(\alph*)]
\item $I=J$.
\item For $\lambda\ge0$ define $\varphi_\lambda(x):=\sup_{y\in S}\{f(y)-\lambda c(x,y)\}$.  There exists a
      dual optimizer of the form $(\lambda,\varphi_\lambda)$ for some $\lambda\ge0$.  Moreover feasible
      $\pi^*\in\Phi_{\mu,\delta}$ and $(\lambda^*,\varphi_{\lambda^*})\in\Lambda_{c,f}$ are primal/dual
      optimizers, $I(\pi^*)=J(\lambda^*,\varphi_{\lambda^*})$, \emph{iff}
      \begin{align}
        f(y)-\lambda^*c(x,y) &= \sup_{z\in S}\{f(z)-\lambda^*c(x,z)\}, \qquad \pi^*\text{-a.s.} \tag{8a}\\
        \lambda^*\Big(\int c(x,y)\dd\pi^*(x,y)-\delta\Big) &= 0. \tag{8b}
      \end{align}
\end{enumerate}
\end{quote}

\paragraph{Remark 1 (the univariate dual, eq.~(9)) --- the form we formalize.}
\begin{equation}\tag{9}
  I \;=\; \inf_{\lambda\ge0}\Big\{\lambda\delta + \E_\mu\Big[\sup_{y\in S}\{f(y)-\lambda c(X,y)\}\Big]\Big\}.
\end{equation}
This follows from (a) together with $J(\lambda,\varphi_\lambda)\le J(\lambda,\varphi)$ for every
$(\lambda,\varphi)\in\Lambda_{c,f}$: among all $\varphi$ dominating $f(y)-\lambda c(x,y)$, the
$c$-transform $\varphi_\lambda$ is the pointwise least.  Measurability of $\varphi_\lambda$ (indeed
$\varphi_\lambda\in mU(S;\R)$) is not immediate and is proved in \S4.

\section{The paper's proof of Theorem 1(a)}

\begin{distillation}
\item \textbf{Weak duality} $I\le J$ is immediate from the definition of $\Lambda_{c,f}$: integrate
      $f(y)\le \varphi(x)+\lambda c(x,y)$ against any $\pi\in\Phi_{\mu,\delta}$.
\item \textbf{Compact $S$, continuous $c$ (Proposition 5).}  Take $X=C_b(S\times S)$, $X^*=M(S\times S)$
      (bounded continuous functions, finite Borel measures; sup and total-variation norms).  Apply the
      \textbf{Fenchel duality theorem} (Luenberger (1997), Ch.~7 Thm 1), exactly as in the standard
      proof of Kantorovich duality on compact spaces (Villani (2003)).  This yields $I=J$ \emph{and}
      existence of a primal optimizer $\pi^*$.
\item \textbf{Compact $S$, lower semicontinuous $c$ (Proposition 6).}  Approximate $c$ from below by
      continuous costs and pass to the limit.
\item \textbf{General Polish $S$.}  The obstruction is that $\Phi_{\mu,\delta}$ is \emph{not tight}
      when $S$ is non-compact --- unlike $\Pi(\mu,\nu)$ in Kantorovich duality.  Exhaust $S$ by compacts
      $S_n$, let $E$ be the convex set of $\pi$ satisfying Proposition 7's conditions, and put
      \[
        T(\lambda,\pi) \;:=\; \lambda\delta + \int \sup_{y\in S_\pi}\{f(y)-\lambda c(x,y)\}\dd\mu(x).
      \]
      $T$ is \emph{concave in $\pi$} for fixed $\lambda$ and convex (indeed affine plus a sup of affines)
      in $\lambda$.  Apply \textbf{Sion's minimax theorem} (Sion (1958)) on $E\times[0,\lambda_{\max}]$
      to interchange $\sup_\pi\inf_\lambda = \inf_\lambda\sup_\pi$, and conclude $J\le I$.  With weak
      duality, $J=I=\max\{I(\pi):\pi\in\Phi_{\mu,\delta}\}$.
\item \textbf{Theorem 1(b)} then identifies the optimizer: the dual optimum is attained at some finite
      $\lambda^*\ge0$ with $\varphi_{\lambda^*}$ the $c$-transform, and (8a)/(8b) are the
      complementary-slackness characterisation.
\end{distillation}

\section{Where our Lean development differs}\label{sec:diff}

\begin{formalizationnotes}
\item \textbf{We formalize eq.~(9), not $I=J$ in the $(\lambda,\varphi)$ form.}  Our
      \texttt{droValue \{m : otCost c m nu <= d\} f = sInf \{lam*d + expect nu phi\_lam : lam > 0\}} is the
      univariate reformulation.  The reduction from $J$ to the $c$-transform (``$\varphi_\lambda$ is
      pointwise least'') is folded into the statement rather than proved.
\item \textbf{Entirely different proof.}  The paper uses \emph{Fenchel duality} on $C_b(S\times S)$ and
      then \emph{Sion's minimax} after a compact exhaustion.  Mathlib has neither in usable form.  Our
      \texttt{ForMathlib.OT.dualValue\_le\_droValue} instead proves the $\ge$ direction by
      \begin{enumerate}[label=(\arabic*),leftmargin=2em]
      \item a \emph{converse Lagrangian bound}: push $\nu$ forward along a \emph{measurable
            $\varepsilon$-argmax} of the $c$-transform (\texttt{exists\_measurable\_eps\_argmax\_of\_separable});
      \item \emph{supergradient existence} for a concave function on $\R$ (\texttt{exists\_nonneg\_multiplier'});
      \item \emph{concavity and monotonicity of the DRO value function} $h(t)=\sup\{\E_\mu[f]: \exists\pi,\ \E_\pi[c]\le t\}$,
            from convexity of the coupling set.
      \end{enumerate}
      Assembled by running (1) at $\lambda=\lambda^*+\eta$, never at $\lambda^*$.  So this is an
      \emph{independent} verification of Theorem 1(a)/eq.~(9), not a transcription of the paper's argument.
\item \textbf{No attainment anywhere.}  Theorem 1(b) asserts a dual optimizer $(\lambda^*,\varphi_{\lambda^*})$
      \emph{exists}, and $\Phi_{\mu,\delta}$ is defined using the fact that $d_c$'s infimum is attained
      (Villani Thm 4.1).  We deliberately use neither: \texttt{otCost} is an \texttt{sInf} with no
      attainment, and an \texttt{sInf} never hands you a minimizer.  Hence the $\lambda^*+\eta$ trick.
      \textbf{Theorem 1(b) --- existence of the dual optimizer and the complementary slackness (8a)/(8b)
      --- is \emph{not} formalized.}  It is the natural next target, and is what the deferred
      worst-case-measure existence theorems (\texttt{GaoKleywegt2023.worstCase\_structure\_cor1},
      \texttt{MohajerinEsfahaniKuhn2018.worstCase\_exists}) still assume.
\item \textbf{Our regularity is stronger; theirs is sharper.}  (A1) asks $c$ lower semicontinuous; we ask
      $c$ \emph{continuous} (\texttt{hcc}), because the $\varepsilon$-argmax selector needs continuity in
      the maximization variable on a separable domain.  (A2) asks $f$ upper semicontinuous and
      $L^1(\mu)$; we ask $f$ continuous and \emph{bounded} ($|f|\le C$).  Both are strictly stronger.
      Weakening $c$ to lsc would require a measurable-selection theorem we deliberately avoided.
\item \textbf{(A1)'s second clause is exactly why Wasserstein needs no Slater.}  We never invoke
      ``$c(x,y)=0\iff x=y$'' by name, but our assembly needs
      \texttt{hne0 : (droValueSet c f nu 0).Nonempty} --- a \emph{zero-cost} feasible plan --- and the
      diagonal coupling supplies it precisely because $c$ vanishes on the diagonal.  With $\delta>0$ the
      radius is then automatically interior to the value function's domain $[0,\infty)$.
      \emph{The entropic objective has no such zero}, which is the whole reason
      \texttt{WangGaoXie2023.strong\_duality} carries a Slater hypothesis and this one does not.  See
      \texttt{WangGaoXie2025\_sinkhorn\_dro\_theorem1.tex}.
\item \textbf{$\lambda\ge0$ vs.\ $\lambda>0$.}  The paper's dual infimum ranges over $\lambda\ge0$; ours
      over $\lambda>0$.  They agree: $\varphi_\lambda\uparrow \sup_y f$ as $\lambda\downarrow0$ by
      monotone convergence, so the $\lambda=0$ value is the limit of the $\lambda>0$ values and cannot
      lower the infimum.  We take $\lambda>0$ because the cost-integrability step needs $\lambda\ne0$.
\end{formalizationnotes}

\TODO{Formalize Theorem 1(b): existence of a finite dual optimizer $\lambda^*$ and the complementary
slackness (8a)/(8b).  This is the missing ingredient behind the deferred worst-case-measure existence
theorems.  Also read \S5 (sufficient conditions for a primal optimizer $\pi^*$) and record which of
those conditions our \texttt{hattain}-style edges correspond to.}

\begin{thebibliography}{99}
\bibitem{ChenGeorgiouPavon2021} Yongxin Chen, Tryphon T. Georgiou, and Michele Pavon. \emph{Stochastic Control Liaisons: Richard Sinkhorn Meets Gaspard Monge on a Schr\"odinger Bridge}. SIAM Review, 63(2):249--313, 2021. doi:10.1137/20M1339982. arXiv:2005.10963.
\bibitem{Leonard2014Survey} Christian L\'eonard. \emph{A Survey of the Schr\"odinger Problem and Some of Its Connections with Optimal Transport}. Discrete and Continuous Dynamical Systems--A, 34(4):1533--1574, 2014. doi:10.3934/dcds.2014.34.1533. arXiv:1308.0215.
\bibitem{Fortet1940} Robert Fortet. \emph{Resolution d'un systeme d'equations de M. Schrodinger}. Journal de mathematiques pures et appliquees, 9e serie, 19(1--4):83--105, 1940.
\bibitem{SinkhornKnopp1967} Richard Sinkhorn and Paul Knopp. \emph{Concerning Nonnegative Matrices and Doubly Stochastic Matrices}. Pacific Journal of Mathematics, 21(2):343--348, 1967.
\bibitem{FranklinLorenz1989} Joel Franklin and Jens Lorenz. \emph{On the Scaling of Multidimensional Matrices}. Linear Algebra and its Applications, 114/115:717--735, 1989. doi:10.1016/0024-3795(89)90490-4.
\bibitem{Carroll2004} Joseph E. Carroll. \emph{Birkhoff's Contraction Coefficient}. Linear Algebra and its Applications, 389:227--234, 2004. doi:10.1016/j.laa.2004.02.039.
\bibitem{EvesonNussbaum1995} Simon P. Eveson and Roger D. Nussbaum. \emph{An Elementary Proof of the Birkhoff--Hopf Theorem}. Mathematical Proceedings of the Cambridge Philosophical Society, 117(1):31--55, 1995.
\bibitem{BlanchetMurthy2019} Jose Blanchet and Karthyek R. A. Murthy. \emph{Quantifying Distributional Model Risk via Optimal Transport}. Mathematics of Operations Research, 44(2):565--600, 2019. doi:10.1287/moor.2018.0936. arXiv:1604.01446.
\bibitem{GaoKleywegt2023} Rui Gao and Anton J. Kleywegt. \emph{Distributionally Robust Stochastic Optimization with Wasserstein Distance}. Mathematics of Operations Research, 48(2):603--655, 2023. doi:10.1287/moor.2022.1275. arXiv:1604.02199.
\bibitem{MohajerinEsfahaniKuhn2018} Peyman Mohajerin Esfahani and Daniel Kuhn. \emph{Data-Driven Distributionally Robust Optimization Using the Wasserstein Metric: Performance Guarantees and Tractable Reformulations}. Mathematical Programming, 171(1--2):115--166, 2018. doi:10.1007/s10107-017-1172-1. arXiv:1505.05116.
\bibitem{WangGaoXie2025} Jie Wang, Rui Gao, and Yao Xie. \emph{Sinkhorn Distributionally Robust Optimization}. Operations Research, 2025. doi:10.1287/opre.2023.0294. arXiv:2109.11926.
\bibitem{Girsanov1960} I. V. Girsanov. \emph{On Transforming a Certain Class of Stochastic Processes by Absolutely Continuous Substitution of Measures}. Theory of Probability and Its Applications, 5(3):285--301, 1960. doi:10.1137/1105027.
\bibitem{CameronMartin1945} R. H. Cameron and W. T. Martin. \emph{Transformations of Wiener Integrals under a General Class of Linear Transformations}. Transactions of the American Mathematical Society, 58(2):184--219, 1945.
\bibitem{Yeh1978} J. Yeh. \emph{Transformation of Conditional Wiener Integrals under Translation and the Cameron--Martin Translation Theorem}. Tohoku Mathematical Journal, 30(4):505--515, 1978.
\bibitem{Kakutani1948} Shizuo Kakutani. \emph{On Equivalence of Infinite Product Measures}. Annals of Mathematics, 49(1):214--224, 1948.
\bibitem{ShiDeBortoliCampbellDoucet2023} Yuyang Shi, Valentin De Bortoli, Andrew Campbell, and Arnaud Doucet. \emph{Diffusion Schr\"odinger Bridge Matching}. Advances in Neural Information Processing Systems, 2023. arXiv:2303.16852.
\bibitem{LiuLipmanNickelKarrerTheodorouChen2024} Guan-Horng Liu, Yaron Lipman, Maximilian Nickel, Brian Karrer, Evangelos A. Theodorou, and Ricky T. Q. Chen. \emph{Generalized Schr\"odinger Bridge Matching}. International Conference on Learning Representations, 2024. arXiv:2310.02233.
\bibitem{DomingoEnrichDrozdzalKarrerChen2025} Carles Domingo-Enrich, Michal Drozdzal, Brian Karrer, and Ricky T. Q. Chen. \emph{Adjoint Matching: Fine-Tuning Flow and Diffusion Generative Models with Memoryless Stochastic Optimal Control}. arXiv preprint arXiv:2409.08861, 2025.

\bibitem{Birkhoff1957} Garrett Birkhoff. \emph{Extensions of Jentzsch's Theorem}. Transactions of the American Mathematical Society, 85(1):219--227, 1957.
\bibitem{Sinkhorn1964} Richard Sinkhorn. \emph{A Relationship Between Arbitrary Positive Matrices and Doubly Stochastic Matrices}. Annals of Mathematical Statistics, 35(2):876--879, 1964.
\bibitem{Sinkhorn1967Prescribed} Richard Sinkhorn. \emph{Diagonal Equivalence to Matrices with Prescribed Row and Column Sums}. American Mathematical Monthly, 74(4):402--405, 1967.
\bibitem{Bauer1965} Friedrich L. Bauer. \emph{An Elementary Proof of the Hopf Inequality for Positive Operators}. Numerische Mathematik, 7:331--337, 1965.
\bibitem{Bushell1973} P. J. Bushell. \emph{Hilbert's Metric and Positive Contraction Mappings in a Banach Space}. Archive for Rational Mechanics and Analysis, 52:330--338, 1973.

\end{thebibliography}

\end{document}
